\section{Methodology}
\label{sec:methodology}

\subsection{Study Design}

Usability testing uses \textbf{moderated task-based sessions}. Each participant works through a fixed set of scenarios (Section~\ref{sec:task_scenarios}) on a prepared device or the participant's device, in a consistent build (e.g., Expo development client or release candidate). The moderator explains the purpose, obtains consent, and encourages a \textbf{think-aloud} protocol unless a participant finds it distracting (note the exception in session notes).

\subsection{Participants}

\begin{itemize}
  \item \textbf{Inclusion:} McMaster students, staff, or faculty (or individuals who match the intended domain, e.g., test accounts with \texttt{@mcmaster.ca}) who are willing to sign consent and have not developed the feature under test.
  \item \textbf{Sample size:} Target a minimum number of sessions (e.g., 5--8) consistent with course guidance; record actual $N$ in Section~\ref{sec:results}.
  \item \textbf{Roles:} Mix riders and drivers where possible, or assign persona per session and document role in the participant log.
\end{itemize}

\subsection{Environment and Materials}

\begin{itemize}
  \item Device: iOS and/or Android; screen recording if permitted by consent.
  \item Network: Stable connection; document if testing against staging or local API.
  \item Test data: Pre-seeded accounts or anonymized personas; reset state between sessions as needed.
  \item \textbf{Ethics:} Brief verbal or written consent covering recording, data handling, and voluntary withdrawal. Store notes and recordings per course or institutional policy.
\end{itemize}

\subsection{Session Outline}

\begin{enumerate}
  \item Welcome and consent (5 min).
  \item Background questionnaire (Appendix~\ref{sec:appendix}, optional short form).
  \item Task scenarios in a fixed order for comparability, or counterbalanced order across participants (document order).
  \item Post-task and post-session questions (Section~\ref{sec:questions_surveys}).
  \item SUS or equivalent survey (Appendix~\ref{sec:appendix}).
  \item Debrief and thanks.
\end{enumerate}

\subsection{Data Recording}

For each task, record: start/end time, success or failure, number and type of assists (minimal prompt vs.\ direct help), and brief qualitative notes (quotes, confusion points). Use a session log template (spreadsheet or form) referenced in Section~\ref{sec:results}.
